\documentclass{beamer}

% +==============================+
% | Vydefinování maker a packagů |
% +==============================+

\usepackage[utf8]{inputenc}
\usepackage[czech]{babel}
\usepackage{graphicx}
\usepackage{tikz}
\usetikzlibrary{shapes.arrows}
\usepackage{MnSymbol,wasysym}
\usepackage{multirow}
\usepackage{xcolor}

\newcommand{\cevroLogo}{
	\begin{tikzpicture}[remember picture,overlay]  
		\node [xshift=-1.6cm,yshift=-0.5cm] at (current page.north east)
		{\includegraphics[width=3cm,height=2cm,keepaspectratio]{images/newcevro}};
	\end{tikzpicture}
}

% +=========================+
% | Vydefinování tikz stylů |
% +=========================+

\tikzset{
    myarrow/.style={
        draw,
        fill=green,
        single arrow,
        minimum height=3.5ex,
        single arrow head extend=1ex
    }
}

\tikzset{
    myarrowred/.style={
        draw,
        fill=red,
        single arrow,
        minimum height=3.5ex,
        single arrow head extend=1ex
    }
}

% +=======================================+
% | Základní informace o bakalářské práci |
% +=======================================+

\title{Svěřenský fond jako nástroj mezigeneračního převodu majetku}
%\subtitle[krátký název prezentace]{dlouhý název prezentace}
%\author{Dominik Bálint}
\author[My name]{Dominik Bálint\\ \footnotesize Vedoucí práce: Mgr. Ing. Tomáš Střeleček, LL.M., Ph.D.\\ \footnotesize Oponent práce: JUDr. Miroslav Sedláček, Ph.D., LL.M.}
\institute{CEVRO Institut}
\date{\today}

% +=============+
% | První sekce |
% +=============+

\section{Prezentace k obhajobě}

% +===============================================+
% | Použitý theme v prezentaci k bakalářské práci |
% +===============================================+

\usetheme{Singapore}

\begin{document}

% +===============+
% | Titulní slide |
% +===============+

\maketitle

% +====================+
% | Začátek prezentace |
% +====================+

\begin{frame}
	\frametitle{Obsah}
	\cevroLogo

\begin{itemize}
	\item Cíl práce%\pause
	\item Metodika%\pause
	\item Výsledky, doporučení%\pause
	\item Závěr%\pause
	\item Zodpovězení otázek
\end{itemize}
	
\end{frame}

% +=============+
% | Druhý slide |
% +=============+

\begin{frame}
	\frametitle{Cíl práce}
	\cevroLogo

\fbox{Cílem práce} jest představit svěřenský fond,\\
jeho použití k mezigeneračnímu převodu majetku,
a zodpovězení \textbf{hypotézy}, zda lze svěřenský fond úspěšně použít k obcházení práva nepominutelných dědiců na povinný díl.

\vspace*{5mm}

\fbox{Vedlejší cíle} : další problémy spojené se svěřenským fondem.

\end{frame}

% +=============+
% | Třetí slide |
% +=============+

\begin{frame}
	\frametitle{Metodika}
	\cevroLogo

\begin{columns}[T]
	\column{0.5\textwidth}
	\textbf{Stanovení cíle}\\
		Z pozitivistického hlediska prozkoumat vztah mezi jednotlivými ustanoveními svěřenského fondu a dědického práva.\\
		---\\
		Účelem dovození odpovědi na vyloženou hypotézu a odhalení problémů.\\
		---\\
		Konstruování normativního výstupu - navržení obecných změn, pokud jich bude třeba.
	\column{0.5\textwidth}
	\textbf{Výzkumné metody}\\
		Obecně teoretické metody.\\
		---\\
		Popisná, kompilační metoda u teoretické části.\\
		---\\
		Analyticko-syntetické metody u praktické části.\\
		---\\
		Zkoumáno (zkoumaná problematika) dvěma optikami ve smyslu interpretace právních norem.\\
		%\vspace*{5.9cm}
\end{columns}

\end{frame}

% +==============+
% | Čtvrtý slide |
% +==============+

\vspace*{-1cm}

\begin{frame}
	\frametitle{Metodika}
	\cevroLogo
	
\includegraphics[height=\textheight,width=\textwidth,keepaspectratio]{images/Zkoumani.png}

\vspace*{0.2cm}

\begin{columns}[T]
\column{0.5\textwidth}
{\small\textbf{Optika 1}}\\
Systematická\\
Logická\\
Teleologická\\
Komparativní (pro ilustraci vzorových řešení)\\
Doktrinální\\
Smysl ve vztahu ke konstrukci (struktuře) svěřenského fondu.\\
\column{0.5\textwidth}
{\small\textbf{Optika 2}	}\\
Formální - povrchní\\
Gramatická, §2\\
"Logická"\\
Dedukce úmyslu zákonodárce\\
\end{columns}

	
\end{frame}

% +============+
% | Pátý slide |
% +============+

\begin{frame}
	\frametitle{Výsledky, doporučení}
	\cevroLogo

\fbox{Cíle jednotlivých částí}\\
\vspace*{5mm}
Teoretická část $\rightarrow$ popisná část (historie, přijetí, právní úprava).\\
Praktická část $\rightarrow$ analytická část (§1400 - §1474 OZ {\color{red} X} §1475 - §1720 OZ).\\
Závěrečná část $\rightarrow$ návrh řešení problematických částí (pokud budou nalezeny - především kolize s dědickým právem).\\
\vspace*{5mm}
\fbox{Výsledek analytické částí}\\
\vspace*{5mm}
Ve vztahu k hypotéze {\color{red} nedosaženo} výsledku $\rightarrow$ hypotéza nezodpovězena, vysloven pouze osobní názor, domněnka ochrany nepominutelných dědiců $\rightarrow$ dovození problematiky, potvrzení problémů, shoda {\color{red} X} konfrontace s literaturou zaměřující se na stejné téma.\\

\end{frame}

% +=============+
% | Šestý slide |
% +=============+

\begin{frame}
	\frametitle{Výsledky, doporučení}
	\cevroLogo

\textbf{Výkladová optika} $\rightarrow$ {\color{red} nedovozená žádná} ochrana.\\
\textbf{Úmyslová optika} $\rightarrow$ {\color{green} potenciálně dovozená} ochrana, kolize s absencí právní osobnosti svěřenského fondu popřípadě její fikce ve vztahu k pozůstalostnímu řízení.\\
-----------------------------------------\\
{\color{red} Proti}: absence právní osobnosti svěřenského fondu, svěřenskému správci nesvědčí vlastnické právo, vyčleněný majetek není součástí pozůstalosti, autonomie vůle, potenciální znevýhodnění svěřenského fondu \textit{mortis causa} oproti svěřenskému fondu \textit{inter vivos}.\\
{\color{green} Pro}: §1452, §311, §1492 OZ, částečně také §17, odstavec 2.\\
Doporučení: novelizace.

\end{frame}

% +=============+
% | Sedmý slide |
% +=============+

\begin{frame}
	\frametitle{Závěr}
	\cevroLogo

\setbeamercolor{block title}{use=structure,fg=white,bg=blue!75!black}
\setbeamercolor{block body}{use=structure,fg=black,bg=white!20!white}
\vspace*{-5mm}
\begin{block}{\small§ 1452}\small{
(1) Každý svěřenský fond musí mít statut. Statut svěřenského\\ fondu vydává zakladatel. Zřizuje-li se svěřenský fond\\ pořízením pro případ smrti, použije se přiměřeně § 311.}
\end{block}
\vspace*{-2mm}
\centering\tikz [baseline=-1ex]{\node [myarrow,rotate=-90] {};}\\
\vspace*{-2mm}
\begin{block}{\small§ 311}\small{
(1) Při založení nadace pořízením pro případ smrti se do nadace vnáší vklad povoláním nadace za dědice nebo nařízením odkazu. V takovém případě nabývá založení nadace účinnosti smrtí zůstavitele.}
\end{block}
\vspace*{-2mm}
\centering\textbf{\Large{\color{green} +}}\\
\vspace*{-2mm}
\begin{block}{\small§ 1492}\small{
Pořízením pro případ smrti nelze zkrátit povinný díl nepominutelného dědice, který se práva na povinný díl nezřekl a nedošlo-li ani k vydědění. Pokud tomu pořízení pro případ smrti odporuje, náleží nepominutelnému dědici povinný díl.}
\end{block}

\end{frame}

% +============+
% | Osmý slide |
% +============+

\begin{frame}
	\frametitle{Závěr}
	\cevroLogo

\begin{center}
Absence kolizní právní úpravy,\\
doktrinálně odlišné pohledy,\\
absence soudních rozhodnutí.\\
\tikz [baseline=-1ex]{\node [myarrowred,rotate=-90] {};}\\
Právní nejistota,\\
není jasné, zda je dědic chráněn,\\
nutnost explicitní úpravy.
\end{center}

\end{frame}


% +==============+
% | Devátý slide |
% +==============+

\begin{frame}
	\frametitle{Zodpovězení otázek}
	\cevroLogo

\begin{table}[]
\begin{tabular}{|l|}
\hline
Otázka 1 Vedoucí práce                                                                                                                                                                                                                       \\ \hline
\begin{tabular}[c]{@{}l@{}}Které otázky problematického vztahu právní úpravy svěřenských\\ fondů a dědického práva si podle autora žádají nejnaléhavěji\\ řešení přijetím novely?\\ Jak by taková novela mohla, podle autora, znít?\end{tabular} \\ \hline
Nejzásadnější změna by měla nastat ve smyslu kolize dědického\\ práva s právem svěřenských fondů.\\
---\\
Další problémy vytyčené v práci (započtení, daňová problematika)\\ jsou oproti této problematice absolutně marginálního charakteru.\\
---\\
Mělo by dojít primárně k řešení problematiky nepominutelných\\ dědiců.                                                                                                                                                                                                                             \\ \hline
\end{tabular}
\end{table}

\end{frame}

% +==============+
% | Desátý slide |
% +==============+

\begin{frame}
	\frametitle{Zodpovězení otázek}
	\cevroLogo

\begin{columns}
	\column{0.5\textwidth}
		Hello everybody
	\column{0.5\textwidth}
		Hello everybody
\end{columns}

\end{frame}

% +=================+
% | Jedenáctý slide |
% +=================+

\begin{frame}
	\frametitle{Zodpovězení otázek}
	\cevroLogo

\begin{table}[]
\begin{tabular}{|l|}
\hline
Otázka 2 Vedoucí práce                                                                                                                                                                                                                       \\ \hline
\begin{tabular}[c]{@{}l@{}}Lze změnit statut svěřenského fondu zřízeného mortis causa,\\ příp. za jakých podmínek? Kdo ho může změnit?\end{tabular} \\ \hline
Zřízení jako založení vs. vznik.\\
---\\
Zřízení (§ 1451) pro případ smrti - založen, zakladatel je živ.\\
---\\
Svěřenský fond neexistuje a proto je možné jej měnit stejně,\\ jako je možné měnit pořízení pro případ smrti.\\
---\\
Dědická dohoda - pouze po dohodě stran - závazek, například\\ novace.                                                                                                                                                                                                                              \\ \hline
\end{tabular}
\end{table}

\end{frame}

% +================+
% | Dvanáctý slide |
% +================+

\begin{frame}
	\frametitle{Zodpovězení otázek}
	\cevroLogo

\begin{table}[]
\begin{tabular}{|l|}
\hline
Otázka 1 Oponent práce                                                                                                                                                                                                                       \\ \hline
\begin{tabular}[c]{@{}l@{}}Nechť autor v rámci ústní obhajoby shrne základní podobu\\ konceptu majetkové struktury svěřenského fondu.\end{tabular} \\ \hline
Hello everybody\\
Hello everybody\\
Hello everybody\\
Hello everybody\\
Hello everybody\\
Hello everybody\\
Hello everybody\\
Hello everybody\\
Hello everybody\\
Hello everybody                                                                                                                                                                                                                              \\ \hline
\end{tabular}
\end{table}

\end{frame}

% +================+
% | Třináctý slide |
% +================+

\begin{frame}
	\frametitle{Zodpovězení otázek}
	\cevroLogo

\begin{table}[]
\begin{tabular}{|l|}
\hline
Otázka 2 Oponent práce                                                                                                                                                                                                                      \\ \hline
\begin{tabular}[c]{@{}l@{}}Autor by mohl dále shrnout některé dílčí odlišnosti québecké\\ a české úpravy svěřenských fondů a tím učinit závěr v dané otázce\\ na str. 20 až 23.\end{tabular} \\ \hline
Hello everybody\\
Hello everybody\\
Hello everybody\\
Hello everybody\\
Hello everybody\\
Hello everybody\\
Hello everybody\\
Hello everybody\\
Hello everybody\\
Hello everybody                                                                                                                                                                                                                              \\ \hline
\end{tabular}
\end{table}

\end{frame}

% +================+
% | Čtrnáctý slide |
% +================+

\begin{frame}

\begin{center}
\Large{Děkuji za pozornost	\smiley{}}
\end{center}

\end{frame}
	
\end{document}
